% Answers to Chapter 16, from lectures/L16/appendix/c_solutions.md; the self-test from
% lectures/L16/README.md.

\facitavsnitt{c16}{Kapitel 16}{Komplexa tal, del II}
\begin{facit}

\svar{1.1} $I \approx 5\,\angle\,0{,}93\,\text{mA}$, med vinkeln $0{,}93\,\text{rad} \approx
53{,}1^{\circ}$.

\svar{1.2} $U = \sqrt{2} + j\sqrt{2} \approx 1{,}41 + j1{,}41\,\text{V}$

\svar{1.3} $Z \approx 0{,}4\,\angle\,{-0{,}14}\,\text{k}\Omega$ (vinkeln i radianer), på
rektangulär form $Z \approx 0{,}4 - j0{,}06\,\text{k}\Omega$; noggrannare
$396 - j57\,\Omega$.

\svar{2.1} \sv{a} $a = 4 + j3$, $b = -3 + j6$, $c = -1 + j10$
\sv{b} Pilar från origo till $(4; 3)$, $(-3; 6)$ och $(-1; 10)$: $a$ i första kvadranten, $b$ och
$c$ i andra.
\sv{c} $\abs{a} = 5$, $\abs{b} = \sqrt{45} \approx 6{,}7$, $\abs{c} = \sqrt{101} \approx 10{,}0$
\sv{d} $b + 2c = -5 + j26$, $\abs{b + 2c} = \sqrt{701} \approx 26{,}5$
\sv{e} $\delta_a \approx 0{,}64\,\text{rad} \approx 36{,}9^{\circ}$,
$\delta_b \approx 2{,}0\,\text{rad} \approx 116{,}6^{\circ}$,
$\delta_c \approx 1{,}7\,\text{rad} \approx 95{,}7^{\circ}$
\sv{f} $d = -8 - j6$

\svar{2.2} \sv{a} En pil från origo till $(-12; 4)$, i andra kvadranten.
\sv{b} $\abs{U} = \sqrt{160} \approx 12{,}6\,\text{V}$, $\delta \approx 2{,}8\,\text{rad} \approx
161{,}6^{\circ}$, alltså $U \approx 12{,}6\,e^{j2{,}8}\,\text{V}$.
\sv{c} $\omega = 200\pi \approx 628\,\text{rad/s}$
\sv{d} $u(t) \approx 12{,}6\,e^{j(200\pi t + 2{,}8)}\,\text{V}$

\svar{3.1} \sv{a} $j$ \sv{b} $-1$ \sv{c} $1$ \sv{d} $-j$ \sv{e} $1$

\svar{3.2} \sv{a} $4$ \sv{b} $2{,}5 + j4{,}33$ \sv{c} $\approx 1{,}41 - j1{,}41$
\sv{d} $-5 + j8{,}66$

\svar{3.3} \sv{a} $3\sqrt{2}\,e^{j\pi/4} \approx 4{,}24\,e^{j\pi/4}$
\sv{b} $2\sqrt{2}\,e^{j3\pi/4} \approx 2{,}83\,e^{j3\pi/4}$
\sv{c} $6\,e^{-j\pi/2}$ \sv{d} $5\,e^{j\pi}$

\svar{3.4} \sv{a} $12\,\angle\,75^{\circ}$
\sv{b} $\frac{4}{3}\,\angle\,{-15^{\circ}} \approx 1{,}33\,\angle\,{-15^{\circ}}$
\sv{c} $16\,\angle\,60^{\circ}$
\sv{d} Beloppen multipliceras och vinklarna adderas; på rektangulär form krävs fyra delprodukter
och hantering av $j^2$.

\svar{3.5} \sv{a} $-1 + j7$
\sv{b} $z_1 \approx 5\,\angle\,53{,}1^{\circ}$, $z_2 = \sqrt{2}\,\angle\,45^{\circ} \approx
1{,}41\,\angle\,45^{\circ}$
\sv{c} $\approx 7{,}07\,\angle\,98{,}1^{\circ}$
\sv{d} Ja: $\abs{-1 + j7} = \sqrt{50} \approx 7{,}07$ och vinkeln är $\approx 98{,}1^{\circ}$.

\svar{3.6} \sv{a} $5\,\angle\,60^{\circ}$ \sv{b} $3\,\angle\,{-90^{\circ}}$
\sv{c} $2\,\angle\,{-90^{\circ}}$
\sv{d} Den minskar med $90^{\circ}$ och beloppet är oförändrat, eftersom
$j = 1\,\angle\,90^{\circ}$: talet vrids ett kvarts varv medurs.

\svar{3.7} \sv{a} $4\,\angle\,60^{\circ}$ \sv{b} $8\,\angle\,90^{\circ} = j8$
\sv{c} $16\,\angle\,120^{\circ}$ \sv{d} $z^n = \abs{z}^n\,\angle\,n\delta$ (de Moivres formel)

\svar{3.8} \sv{a} $10\,\angle\,\frac{\pi}{6} = 10\,\angle\,30^{\circ}$
\sv{b} $4\,\angle\,{-\frac{\pi}{2}} = 4\,\angle\,{-90^{\circ}}$
\sv{c} $6\,\angle\,0^{\circ}$
\sv{d} Signalerna måste ha samma vinkelhastighet $\omega$.

\svar{3.9} \sv{a} $u(t) = 8\sin\left(100\pi t + \frac{\pi}{4}\right)\,\text{V}$
\sv{b} $u(t) = 5\sin\left(100\pi t - \frac{\pi}{2}\right)\,\text{V}$
\sv{c} $i(t) \approx 5\sin(100\pi t + 0{,}927)\,\text{mA}$
\sv{d} $i(t) \approx 2{,}83\sin(100\pi t - 2{,}36)\,\text{mA}$

\svar{3.10} \sv{a} $U_1 = 6 + j0\,\text{V}$, $U_2 = 0 + j8\,\text{V}$ \sv{b} $6 + j8\,\text{V}$
\sv{c} $10\,\angle\,53{,}1^{\circ}\,\text{V}$
\sv{d} $u(t) \approx 10\sin(100\pi t + 0{,}927)\,\text{V}$

\svar{3.11} \sv{a} $Z = 5\,\angle\,45^{\circ}\,\Omega$ \sv{b} $Z \approx 3{,}54 + j3{,}54\,\Omega$
\sv{c} Induktiv: imaginärdelen är positiv.
\sv{d} $R \approx 3{,}54\,\Omega$ och $X \approx 3{,}54\,\Omega$

\svar{3.12} \sv{a} $Z_L = j20\,\Omega$ \sv{b} $Z_C = -j20\,\Omega$ \sv{c} $Z = 15\,\Omega$
\sv{d} Resonans: reaktanserna tar ut varandra, $\omega L = 1/(\omega C)$, och impedansen blir
rent resistiv.

\svar{3.13} \sv{a} $a = 5 + j12$, $b = -8 + j6$ \sv{b} $\abs{a} = 13$, $\abs{b} = 10$
\sv{c} $a + b = -3 + j18$, $\abs{a + b} = \sqrt{333} \approx 18{,}25$
\sv{d} $\delta_a \approx 67{,}4^{\circ}$, $\delta_b \approx 143{,}1^{\circ}$

\svar{3.14} \sv{a} $jz = -4 + j3$ \sv{b} Lika stora: $\abs{z} = \abs{jz} = 5$.
\sv{c} $53{,}1^{\circ}$ och $143{,}1^{\circ}$; skillnaden är $90^{\circ}$.
\sv{d} Den vrider talet $90^{\circ}$ moturs utan att ändra dess längd.

\svar{3.15} \sv{a} Sant: $\abs{e^{j\delta}} = \sqrt{\cos^2\delta + \sin^2\delta} = 1$.
\sv{b} Falskt: beloppen multipliceras, $\abs{z_1z_2} = \abs{z_1}\abs{z_2}$.
\sv{c} Sant: $z_1/z_2 = (\abs{z_1}/\abs{z_2})\,\angle\,(\delta_1 - \delta_2)$.
\sv{d} Falskt: en fasor bär bara amplitud och fas; $\omega$ anges separat.
\sv{e} Sant: det är Eulers identitet, $e^{j\pi} = -1$.
\sv{f} Falskt: multiplikation är enklast på polär form eller Eulerform.

\svar[Testa dig själv]{T} \sv{1} $Z = 5\sqrt{3} + j5 \approx 8{,}66 + j5$
\sv{2} $e^{j\pi} = -1$: talet $1$ vrids ett halvt varv ($180^{\circ}$) längs enhetscirkeln och
hamnar i $-1$.

\end{facit}
